\section{System Overview}
\label{sec:overview}

\subsection{Design principles}
Three ideas organize the design.
\textbf{(P1) Attestation turns distrust into unreliability.} Everything that touches tenant secrets or makes scheduling
decisions runs inside the attested appliance, so the rest of the system only has to tolerate omission and delay
(Section~\ref{sec:trust}).
\textbf{(P2) Trust is a managed resource.} Attestation, key provisioning and warm state have costs, lifetimes and an
exposure price, and they are planned on a slower timescale than invocation placement (Section~\ref{sec:sched}).
\textbf{(P3) Liveness is a hint; progress is evidence.} Signals from untrusted hosts, including heartbeats, drain notices
and self-reported health, are never sufficient. Failures are detected where work is observed by a trusted or attested party
(Section~\ref{sec:ft}).

\subsection{Architecture and roles}
Fig.~\ref{fig:arch} shows the components and Table~\ref{tab:roles} states what each is trusted for. Four kinds of
parties take part.

\begin{figure*}[t]
\centering
\resizebox{\textwidth}{!}{% Architecture figure (Fig. 1). Native width 139 mm (acmart single column). y runs downwards.
% Rows: clients and legend; provider control plane; volunteer cluster; tenant domain.
\begin{tikzpicture}[x=1mm,y=1mm]
  \tikzset{box/.append style={font=\footnotesize}}
  % ---------- clients and legend ----------
  \node[box, minimum width=20mm, minimum height=7mm] (client) at (52,-3.5) {Clients};
  \begin{scope}[font=\tinyx, every node/.style={inner sep=1pt}]
    \node[prov, minimum size=3mm] (l1) at (75,-3.5)  {};  \node[anchor=west] at (l1.east) {provider};
    \node[att,  minimum size=3mm] (l2) at (92,-3.5)  {};  \node[anchor=west] at (l2.east) {attested code};
    \node[ten,  minimum size=3mm] (l3) at (114,-3.5) {};  \node[anchor=west] at (l3.east) {tenant};
    \node[untr, minimum size=3mm] (l4) at (128,-3.5) {};  \node[anchor=west] at (l4.east) {untrusted};
  \end{scope}
  % ---------- provider control plane ----------
  \node[zone, fill=blue!4, minimum width=139mm, minimum height=29mm, anchor=north west] (prov) at (0,-10) {};
  \node[font=\scriptsize\bfseries, anchor=north west] at (2,-10.8) {Provider control plane};
  \node[font=\tinyx, anchor=north east] at (137,-11.2) {availability and accounting, no tenant keys};
  \node[prov, minimum width=31mm, minimum height=15mm, font=\scriptsize] (ver)  at (17,-29)  {\textbf{Verifier, registry}\\[-1pt]\tinyx session certificates\\[-2pt]\tinyx encrypted bundles};
  \node[prov, minimum width=31mm, minimum height=15mm, font=\scriptsize] (rtr)  at (52,-29)  {\textbf{Routers}\\[-1pt]\tinyx pending copies\\[-2pt]\tinyx progress watchdogs\\[-2pt]\tinyx receipts};
  \node[prov, minimum width=31mm, minimum height=15mm, font=\scriptsize] (brk)  at (87,-29)  {\textbf{Broker}\\[-1pt]\tinyx one topic set\\[-2pt]\tinyx per cluster};
  \node[prov, minimum width=31mm, minimum height=15mm, font=\scriptsize] (orch) at (122,-29) {\textbf{Orchestrator}\\[-1pt]\tinyx clusters, epochs\\[-2pt]\tinyx database};
  % ---------- volunteer cluster ----------
  \node[zone, fill=gray!5, minimum width=139mm, minimum height=44mm, anchor=north west] (vol) at (0,-48) {};
  \node[font=\scriptsize\bfseries, anchor=north west] at (2,-48.8) {Volunteer cluster: tens of machines};
  \node[att, minimum width=84mm, minimum height=21mm, font=\scriptsize] (w) at (45,-69)
       {\textbf{Worker (attested appliance)}\\[2pt]{\tinyx\begin{tabular}{@{}l@{\hspace{4mm}}l@{}}agent and TPM quote & keys in memory only\\ Firecracker microVMs & sealed, encrypted cache\end{tabular}}};
  \node[att, minimum width=38mm, minimum height=21mm, font=\scriptsize] (sch) at (115,-69)
       {\textbf{Worker with}\\[-1pt]\textbf{scheduler role}\\[-1pt]\tinyx planner and placer\\[-2pt]\tinyx soft state, epoch $e$};
  \node[untr, minimum width=73mm, minimum height=6mm, font=\tinyx] (own) at (99,-87) {Owner controls hardware, power, network: untrusted};
  % ---------- tenant domain ----------
  \node[zone, fill=orange!6, minimum width=66mm, minimum height=21mm, anchor=north west] (tenant) at (3,-96) {};
  \node[font=\scriptsize\bfseries, anchor=north west] at (5,-96.8) {Tenant domain};
  \node[ten, minimum width=58mm, minimum height=8mm, font=\scriptsize] (kbs) at (36,-110) {\textbf{Key service}: secrets to attested hosts only};
  % ---------- numbered flows (one invocation); dashed = control traffic off the hot path ----------
  \draw[msg] (client.south) -- node[lbl, left=0.5mm] {\circled{1}} (rtr.north);
  \draw[msg] (rtr.east) -- node[lbl, above=2.5mm] {\circled{2}} (brk.west);
  \draw[msg, <->] (99,-36.5) -- node[lbl, right=0.5mm, pos=0.3] {\circled{2}\circled{3}} (99,-58.5);
  \draw[msg] (82,-36.5) -- node[lbl, right=0.5mm, pos=0.3] {\circled{3}} (82,-58.5);
  \draw[<->, thick] (52,-36.5) -- node[lbl, right=0.5mm, align=left, pos=0.3] {\circled{4} input, result, receipt} (52,-58.5);
  \draw[msg, <->] (11,-36.5) -- node[lbl, left=0.3mm, pos=0.3] {\circled{5}} (11,-58.5);
  \draw[msg, <-] (23,-58.5) -- node[lbl, right=0.5mm, pos=0.3] {\circled{7}} (23,-36.5);
  \draw[msg, dashed, <->] (brk.east) -- node[lbl, above=2.5mm] {\circled{8}} (orch.west);
  \draw[msg, <->] (36,-79.5) -- node[lbl, right=0.5mm] {\circled{6} key request} (36,-96);
  \node[zone, dashed, fill=gray!4, minimum width=58mm, minimum height=21mm, anchor=north west, font=\tinyx, align=center] (other) at (77,-96) {Other clusters have the same structure\\and their own topics on the same broker};
\end{tikzpicture}
}
\caption{\voss{} architecture. Shading marks who is trusted for what: blue, the provider-operated control plane (availability and
accounting, never tenant secrets); green, attested code running on owner-controlled hardware; orange, the tenant-operated key
broker; grey, the owner's hardware, which is untrusted. Numbered flows follow one invocation (Fig.~\ref{fig:seq}); dashed arrows are
control traffic off the hot path.}
\label{fig:arch}
\end{figure*}

\textbf{Provider control plane.} A provider-operated tier is assumed reliable and is trusted for \emph{availability and
accounting}, but never for tenant confidentiality. \emph{Regional routers} receive invocations, choose a cluster, keep a pending
copy until a worker accepts, observe progress, and write the receipts that are the ground truth of work. The
\emph{orchestrator} partitions workers into geographic clusters of tens of hosts (which bounds the planner's problem size and
intra-cluster latency), appoints one scheduler per cluster, and rebalances clusters. The \emph{verifier and registrar} appraises
attestation evidence and issues session certificates. An \emph{operational database} private to this tier stores identities,
sessions, cluster epochs, checkpoints and receipts, and an \emph{artifact registry} serves encrypted code bundles and root
file systems.

\textbf{Tenants.} Each tenant runs a \emph{key broker service} (KBS) that releases the tenant's secrets only to hosts with an
acceptable attestation result (Section~\ref{sec:provisioning}). Neither the scheduler nor the control plane ever sees a tenant key.

\textbf{Volunteer hosts.} A volunteer machine reboots into the \emph{attested appliance}, a locked, measured operating-system
image containing the worker agent, the Firecracker microVM monitor \cite{agache2020firecracker} and the scheduler code. The
\emph{worker} runs each invocation in its own microVM and keeps a content-addressed, encrypted cache of code and snapshots. The
\emph{scheduler} is a role: the same image runs on every worker, so any worker can be promoted without booting anything, and one
worker per cluster is appointed to perform the scheduling functions of Section~\ref{sec:sched}. The host itself (power, network,
firmware settings) remains under the owner's control, which is why the scheduler's \emph{host} can still omit or delay even though
its \emph{code} cannot misbehave.

\begin{table}[t]
\caption{Roles, location, and what each is trusted for.}
\label{tab:roles}
\centering\scriptsize
\setlength{\tabcolsep}{2.5pt}
\begin{tabular}{@{}p{1.55cm}p{1.45cm}p{2.15cm}p{2.7cm}@{}}
\toprule
Role & Runs on & Trusted for & Not trusted for \\
\midrule
Router & provider & availability, accounting & tenant secrets (never receives them) \\
Orchestrator, verifier, DB, registry & provider & availability, accounting, appraisal & tenant secrets \\
Tenant KBS & tenant & releasing secrets per policy & availability of volunteer hosts \\
Scheduler & role on an attested worker & correct placement (attested code) & availability: its host may omit or delay \\
Worker & volunteer host, attested appliance & executing code faithfully & availability; anything outside the attested stack \\
\bottomrule
\end{tabular}
\end{table}

\subsection{Life of an invocation}
Fig.~\ref{fig:seq} follows one synchronous invocation. A router picks a cluster by round-trip time, health and provisioned
capacity for the tenant, and sends only \emph{metadata} to that cluster's scheduler. The scheduler chooses a worker (and possibly a hedge),
dispatches it, and tells the router who was assigned; the router streams the payload to the worker over a TLS session terminated inside
the attested appliance. The scheduler therefore never touches payloads, results or tenant keys. If the worker is not yet provisioned for
the tenant it requests the key from the tenant's KBS, then climbs the start-up tiers (Section~\ref{sec:lifecycle}) and runs the function.
The result returns directly to the router together with a receipt signed by the worker's session key (Section~\ref{sec:receipts}); the
latency visible to the client ends there. Every term of Eq.~\eqref{eq:latency} is logged separately per invocation, which is what lets
Section~\ref{sec:results} attribute latency to its causes.

\begin{figure*}[t]
\centering
\begin{tikzpicture}[x=1mm,y=1mm]
  \def\seqh{86}
  \seqactor{0}{cl}{Client}
  \seqactor{34}{rt}{Router}
  \seqactor{68}{sc}{Scheduler}
  \seqactor{102}{wk}{Worker}
  \seqactor{136}{kb}{Tenant KBS}
  \seqmsg{-9}{cl}{rt}{\circled{1} invoke(fn, payload)}{msg}
  \seqnote{34}{-15.5}{choose cluster by RTT, health,\\provisioned capacity; keep copy}
  \seqmsg{-22}{rt}{sc}{\circled{2} ENQUEUE (metadata only)}{msg}
  \seqnote{68}{-28.5}{candidates $=$ provisioned workers $\cup$ cold path;\\cost model; place, maybe hedge}
  \seqmsg{-35}{sc}{wk}{\circled{3} DISPATCH(epoch, token)}{msg}
  \seqmsg{-39}{sc}{rt}{ASSIGNED}{msg}
  \seqmsg{-44}{rt}{wk}{\circled{4} INPUT (attested TLS)}{msg}
  \seqmsg{-49}{wk}{rt}{ACCEPTED}{rsp}
  % alt: trust-cold path
  \draw[dashed, gray] (88,-52.5) rectangle (146,-66.5);
  \node[font=\tinyx, anchor=north west, inner sep=1pt] at (88,-52.5) {alt: worker not yet provisioned for tenant};
  \seqmsg{-58}{wk}{kb}{\circled{5} KEY\_REQUEST(cert, evidence)}{msg}
  \seqmsg{-63}{kb}{wk}{key wrapped to session key}{rsp}
  \seqnote{102}{-71.5}{start-up tier: warm sandbox $\succ$ restore snapshot\\$\succ$ boot from disk $\succ$ fetch and decrypt}
  \seqmsg{-79}{wk}{rt}{\circled{6} RESULT $+$ attested receipt}{msg}
  \seqmsg{-83}{wk}{sc}{OUTCOME (timings, tier)}{rsp}
  % latency terms
  \draw[decorate, decoration={brace, amplitude=2pt, mirror}] (152,-12) -- (152,-34) node[midway, right=1mm, font=\tinyx, align=left] {$T_{\text{route}}$\\$T_{\text{decide}}$};
  \draw[decorate, decoration={brace, amplitude=2pt, mirror}] (152,-52.5) -- (152,-66.5) node[midway, right=1mm, font=\tinyx] {$T_{\text{prov}}$};
  \draw[decorate, decoration={brace, amplitude=2pt, mirror}] (152,-68) -- (152,-79) node[midway, right=1mm, font=\tinyx, align=left] {$T_{\text{start}}$\\$T_{\text{exec}}$};
\end{tikzpicture}

\caption{Life of a synchronous invocation. The scheduler handles only metadata. The shaded alternative is the path through an unprovisioned worker, which the planner tries to keep off the critical path.
Braces map messages to the terms of Eq.~\eqref{eq:latency} ($T_{\text{queue}}$ and $T_{\text{net}}$ occur between steps 4 and 6).}
\label{fig:seq}
\end{figure*}

\subsection{Sessions and start-up tiers}
\label{sec:lifecycle}
A worker alternates between OFF and a \emph{session} (Fig.~\ref{fig:lifecycle}a). It becomes schedulable only after booting
and being attested, which on our calibration machine takes \CalAttested\,s. A session ends abruptly (reboot, power loss, crash) or
gracefully (the owner announces a drain with a deadline). Session end discards every secret and warm sandbox, but disk
caches remain as ciphertext and are valid again after the next attestation.

Within a session a function on a host sits at one of the start-up tiers of Section~\ref{sec:background}
(Fig.~\ref{fig:lifecycle}b), and a separate flag $K[t,j]$ records whether the host holds tenant $t$'s key. Start-up costs depend on the tier
(warm microVM boot \CalFcBootWarmMs\,ms with a warm page cache, \CalFcBootColdS\,s with a cold one; snapshot restore
\CalRestoreWarmMs\,ms warm, \CalRestoreColdS\,s cold; Fig.~\ref{fig:cost}) and on whether $K[t,j]=1$. These two facts, a per-session trust step and a
per-tenant key flag, are what the generic lifecycle engine in our simulator interprets from data for every substrate (Section~\ref{sec:substrates}).

\begin{figure*}[t]
\centering
\resizebox{\textwidth}{!}{\begin{tikzpicture}[x=1mm,y=1mm]
  % ---------- (a) volunteer session ----------
  \node[font=\scriptsize\bfseries, anchor=west] at (0,6) {(a) Volunteer session};
  \node[state] (off) at (8,-24) {OFF};
  \node[state] (boot) at (36,-24) {BOOTING};
  \node[state, minimum width=26mm, minimum height=14mm, align=center] (att) at (72,-24) {ATTESTED\\[-1pt]\tinyx schedulable;\\[-2pt]\tinyx provisioned per tenant};
  \node[state] (drain) at (72,-1) {DRAINING};
  \draw[msg] (off.east) to[bend left=18] node[lbl,above]{owner starts} (boot.west);
  \draw[msg, gray] (boot.south west) to[bend left=18] node[lbl,below]{abort} (off.south east);
  \draw[msg] (boot) -- node[lbl,above,align=center]{quote verified} (att);
  \draw[msg] (att.north) -- node[lbl,right,align=left]{drain request\\(deadline; a hint)} (drain.south);
  \draw[msg] (drain.west) -| node[lbl,pos=0.25,above]{deadline or idle} (off.north);
  \draw[msg, red!70!black] (att.south) -- ++(0,-8) -| node[lbl,pos=0.2,below]{reboot, power loss, crash (abrupt)} (off.south);
  \node[font=\tinyx, align=left, anchor=north west, text width=82mm] at (0,-45) {Session end loses the key set $K$ and all warm VMs; disk caches survive as ciphertext and are valid after the next attestation. The session certificate expires at $\min(t_{\text{end}},\,t_{\text{last quote}}+T_{\text{valid}})$ and is extended by periodic re-quotes.};
  % ---------- (b) start-up tiers ----------
  \begin{scope}[shift={(96,0)}]
  \node[font=\scriptsize\bfseries, anchor=west] at (0,6) {(b) Start-up tiers of one function on one host};
  \node[state, minimum width=26mm] (l0) at (36,-46) {L0 nothing cached};
  \node[state, minimum width=26mm] (l1) at (36,-34) {L1 code, rootfs on disk};
  \node[state, minimum width=26mm] (l2) at (36,-22) {L2 snapshot on disk};
  \node[state, minimum width=26mm] (l3) at (36,-10) {L3 warm VM in memory};
  \node[state, minimum width=26mm] (l4) at (36,2) {L4 busy};
  \draw[msg] (l0) -- node[lbl,right]{fetch, verify} (l1);
  \draw[msg] (l1) -- node[lbl,right]{snapshot} (l2);
  \draw[msg] (l2) -- node[lbl,right]{restore} (l3);
  \draw[msg] (l3) -- node[lbl,right]{dispatch} (l4);
  \draw[msg] (l1.west) -- ++(-5,0) |- node[lbl,pos=0.25,left,align=right]{boot, init} (l3.west);
  \draw[decorate, decoration={brace, amplitude=2.5pt}] ([xshift=11mm]l3.north east |- 0,-3) -- ([xshift=11mm]l1.south east |- 0,-37);
  \node[ten, anchor=west, minimum width=24mm, align=left] (k) at (63,-20) {\textbf{Provisioned $K[t,j]$}\\[-1pt]\tinyx tenant key in memory;\\[-2pt]\tinyx needed to decrypt\\[-2pt]\tinyx L1--L3, lost with\\[-2pt]\tinyx the session};
  \end{scope}
\end{tikzpicture}
}
\caption{(a) A volunteer session. Trust is established once per session (booting, attestation) and lost when it ends.
(b) Start-up tiers of one function on one host. Moving up costs time and memory or disk; the tenant key flag $K[t,j]$ is
orthogonal and is required to decrypt the cached artifacts.}
\label{fig:lifecycle}
\end{figure*}

\subsection{State management and communication rules}
Six rules keep the hot path fast and failures recoverable.
(i)~\emph{All cross-node interaction is a typed message}; actors never call each other. Workers sit behind NAT, so each opens outbound
persistent TLS streams to its scheduler, the verifier, and routers, and pushes ride those streams.
(ii)~\emph{The operational database is private to the control plane}, and workers and schedulers reach it only through
messages.
(iii)~\emph{The invocation hot path never touches the database}; a cloud round trip would land in latency.
(iv)~\emph{Live state stays where it is born}: queues, caches and keys in workers, undispatched work in routers, cost matrices in the
scheduler. After a scheduler change it is recovered by asking the workers, not by reading it back from the database
(Section~\ref{sec:succession}).
(v)~\emph{Routers write the ground truth of work}, as attested receipts.
(vi)~\emph{Every channel terminates inside an attested party}: the appliance, the control plane, or the KBS. The session certificate
issued after attestation is the credential for all of them.

\subsection{Component contracts}
\label{sec:contracts}
Table~\ref{tab:contracts} states for each decision-making component what it consumes, what it produces, what it must respect and what it is allowed to know, so that every behaviour in the following sections can be traced to an input and an owner. Parameters are given in Table~\ref{tab:defaults} (scheduling), Table~\ref{tab:ftparams} (fault tolerance) and the calibration bundle (costs).

\begin{table*}[t]
\caption{Contracts of the decision-making components. ``Knows'' lists the information class (Table~\ref{tab:info}).}
\label{tab:contracts}
\centering\scriptsize
\setlength{\tabcolsep}{3pt}
\begin{tabular}{@{}p{1.9cm}p{4.0cm}p{3.6cm}p{4.2cm}p{3.3cm}@{}}
\toprule
Component & Input & Output & Constraints & Knows \\
\midrule
Router (regional) & invocation, cluster summaries each 1\,s & cluster choice, spill, receipt rows & dispatch watchdog 1\,s; sync spill after 1.5\,s & O, measured cold prices \\
Orchestrator & scheduler keep-alives and summaries, worker joins and leaves, stall reports & cluster membership, scheduler appointments (epoch CAS), revocations & single writer per datum; epoch fencing; reliable by assumption & O \\
Scheduler: planner & view, forecasts, tenant policy & plan $\{y,s,w,c\}$, actions & robust coverage, budgets, domain caps, assurance floor & O, H, F, C, P \\
Scheduler: placer & invocation metadata, view, plan & host, hedge, reject or defer & feasibility, hedge budget, admission rule & O, H, C, P \\
Verifier & quote, event log, EK chain & session certificate (assurance, expiry) or rejection & reference measurements, validity period & measured-boot evidence \\
Tenant KBS & certificate, policy & wrapped key or denial & tenant policy only & tenant policy \\
Worker agent (attested) & dispatch, plan actions, input & result and receipt, heartbeat & envelope declared by owner; epoch fencing & local state \\
\bottomrule
\end{tabular}
\end{table*}
